\documentclass[11pt]{article}
\usepackage[a4paper,margin=1in]{geometry}
\usepackage{amsmath,amssymb}
\usepackage{booktabs,longtable}
\usepackage{graphicx}
\usepackage[hyphens]{url}
\usepackage[colorlinks=true,linkcolor=black,urlcolor=blue,citecolor=blue]{hyperref}
\setlength{\parskip}{0.55em}\setlength{\parindent}{0pt}
\providecommand{\tightlist}{\setlength{\itemsep}{0.2em}\setlength{\parskip}{0pt}}

\title{Does a one-pool water-balance model improve forecasts of Edwards Aquifer head? A scored test at J-17}
\author{\textbf{Amin Abaee}\\ Independent Researcher\\[2pt] ORCID: \href{https://orcid.org/0000-0002-0019-1842}{0000-0002-0019-1842}\\ Email: \href{mailto:amin_abaee@ut.ac.ir}{amin_abaee@ut.ac.ir}}
\date{September 18, 2026}
\begin{document}
\maketitle

\section{Does a One-Pool Water-Balance Model Improve Forecasts of
Edwards Aquifer Head? A Scored Test at
J-17}\label{does-a-one-pool-water-balance-model-improve-forecasts-of-edwards-aquifer-head-a-scored-test-at-j-17}

\textbf{Prepared in the format of Water Resources Research.}

\subsection{Highlights}\label{highlights}

\begin{itemize}
\tightlist
\item
  One-pool persistence-recharge balance loses to naive persistence at
  one year: 14.70 ft versus 13.23 ft
\item
  Univariate AR(1) gains 0.39 ft at h = 1 --- a margin within noise (MAE
  tie; five-year loss; bootstrap interval covering zero)
\item
  Climatological-flux affine map: best one-step score (12.28 ft),
  declined by a protocol class clause
\item
  Given realized fluxes the same balance nowcasts (7.55 ft); climate
  modules gain at most 0.13 ft; none retained
\item
  At five years, training-mean climatology beats persistence robustly:
  16.80 versus 21.11 ft
\end{itemize}

\subsection{Abstract}\label{abstract}

\textbf{Problem.} Operational groundwater management frequently relies
on multi-month to multi-year index-well water-level forecasts to trigger
drought stages, enforce critical-period pumping reductions, and maintain
environmental springflows. While modern research has prioritized
increasingly complex data-driven and machine-learning architectures,
simple, lumped process-based water-balance models remain widely used in
conceptual planning. Crucially, such models are rarely evaluated against
naive statistical baselines under a formally frozen model-retention
protocol.

\textbf{Approach.} Using the 90-year annual record (1934--2023) of water
levels at index well J-17 in the San Antonio Pool of the Edwards
(Balcones Fault Zone) Aquifer, Texas, we evaluate a forward-ordered
ladder of discrete-time models using both fixed-window and
rolling-origin out-of-sample root mean squared error (RMSE). The ladder
comprises naive persistence, training-mean climatology, a univariate
autoregression (AR(1)), a one-pool stock-flow water balance, residual
autoregressive and delay formulations, and climate-informed recharge
modules. A causal module is retained if and only if it reduces RMSE
relative to both last-value persistence and the next-simpler causal
model under an evaluation protocol locked prior to scoring.

\textbf{Results.} The one-pool stock-flow model that persists
preceding-year recharge fails to beat naive persistence at the one-year
horizon (\(14.70\text{ ft}\) vs.~\(13.23\text{ ft}\)), because annual
recharge is essentially temporally uncorrelated
(\(\operatorname{corr}(R_t, R_{t-1}) = 0.17\)), whereas annual head
increments strongly reflect contemporaneous recharge
(\(\operatorname{corr}(\Delta H_t, R_t) = 0.74\)). The univariate AR(1)
achieves a marginal one-year RMSE reduction of \(0.39\text{ ft}\)
(\(12.84\text{ ft}\)), but this advantage is statistically
indistinguishable from zero (mean absolute error is tied at
\(10.7\text{ ft}\), persistence wins at five years, and the 95\%
moving-block bootstrap confidence interval spans zero:
{[}\(-1.51, +0.71\){]} ft). A water-balance model driven by
climatological mean fluxes achieves the lowest one-step error
(\(12.28\text{ ft}\)), but represents an affine AR(1) mapping and is
declined on class grounds by a pre-specified protocol clause. When
provided with realized future fluxes, the one-pool balance achieves a
nowcast RMSE of \(7.55\text{ ft}\). Pre-season climate predictors (ENSO
indices and division precipitation) improve one-year head RMSE by at
most \(0.13\text{ ft}\), failing to warrant retention. At a five-year
horizon, training-mean climatology decisively outperforms persistence
(\(16.80\text{ ft}\) vs.~\(21.11\text{ ft}\)).

\textbf{Implications.} For annual water-level forecasting at J-17, naive
persistence and simple univariate autoregression represent exceptionally
demanding benchmarks. Given realized fluxes, the lumped water balance
functions effectively as an annual nowcast rather than an operational
forecast. For multi-year strategic planning, long-term historical
climatology provides a substantially more reliable baseline than dynamic
water balances forced by persisted recharge.

\textbf{Keywords:} Edwards Aquifer, J-17 index well, groundwater level
forecasting, water balance, forecast evaluation, persistence benchmark,
prediction skill

\textbf{Article Impact Statement:} Deliberately simple process-based
water-balance models must justify their complexity against naive
baselines; at J-17, annual recharge is near-white, causing stock-flow
forecasts to lose to persistence unless future fluxes are already known.

\begin{center}\rule{0.5\linewidth}{0.5pt}\end{center}

\subsection{1. Introduction}\label{introduction}

Forecasting groundwater elevations at designated index wells is a
foundational requirement for operational water resources management. In
heavily utilized regional aquifer systems, statutory drought-stage
declarations, mandatory irrigation cutbacks, municipal conservation
rules, and environmental springflow protections are triggered directly
by threshold water-level crossings at representative monitoring
locations. The hydrogeological literature has responded with an
extensive catalog of predictive methodologies, spanning analytical water
balances, lumped conceptual reservoirs, spatially distributed numerical
models, and data-driven formulations, including artificial neural
networks (ANNs) and wavelet-conjunction models (Daliakopoulos et al.,
2005; Adamowski \& Chan, 2011).

Concurrently, the groundwater discipline has begun adopting standardized
benchmark cultures. For instance, the GEMS-GER initiative established a
large-scale machine-learning benchmark dataset spanning 32 years of
weekly observations across 3,207 monitoring wells in Germany, pairing
observational records with three reference models of increasing
sophistication (Ohmer et al., 2026). Similarly, comprehensive
benchmarking across karst catchments has evaluated up to nine
machine-learning and deep-learning architectures against empirical
hydrological indices (Zhu et al., 2026).

When interpreting an index-well record, three distinct hydrogeological
entities must be conceptually separated: 1. \textbf{The Observation Head
(\(H\)):} A local piezometric measurement taken at an access point (the
index well). 2. \textbf{The Aquifer Storage (\(S\)):} The regional
volume of groundwater stored in the porous and karstified rock matrix,
which constitutes the actual resource. 3. \textbf{The Boundary Fluxes
(\(R, P, Q\)):} The recharge pulses, anthropogenic well withdrawals, and
natural spring outflows that drive storage changes.

A lumped, one-pool water-balance model seeks to link these domains by
approximating the dynamic mass balance through an increment structure:
annual head change is modeled as a weighted linear combination of
surface boundary fluxes and a linear regional drainage term. We use the
term \emph{water balance} strictly in this increment-structure sense.
The formulation stores historical spring discharge but deliberately
excludes it as an exogenous predictor in forward forecasting equations,
while enforcing an empirical water-level clipping window ({[}610, 710{]}
ft) as a proxy for physical storage limits. Crucially, this
representation does not constitute a fully closed, conservative physical
mass balance; it is a discrete-time lumped approximation whose
operational forecast utility must be established empirically.

Whether structural elaboration improves genuine \emph{out-of-sample}
predictive skill---as opposed to merely optimizing in-sample
mathematical fit---is a distinct and testable empirical question. In the
broader forecasting discipline, the necessity of rigorous baseline
benchmarks has been repeatedly confirmed. Across the 100,000 empirical
time series analyzed in the landmark M4 competition, complex predictive
algorithms failed to uniformly outperform basic statistical benchmarks
(Makridakis et al., 2020).

Consequently, the minimal performance threshold for any proposed
predictive module is whether it can beat naive last-value persistence
(the forecast of zero change) and long-term training-mean climatology.
Contemporary groundwater benchmarking frameworks typically compare
advanced algorithmic families against one another without subjecting
each intermediate physical mechanism to a pre-registered model-retention
gate scored against naive baselines.

This study provides that missing test. We deploy a formal model-ablation
protocol to evaluate an incrementally structured ladder of lumped
water-balance models against naive baselines, retaining mathematical
complexity if and only if it yields demonstrable out-of-sample error
reductions on the predictand itself.

The San Antonio Pool of the Edwards (Balcones Fault Zone) Aquifer,
Texas, monitored at index well J-17, serves as an ideal test bed (Figure
1). Spanning from 1934 to 2023, the J-17 daily water-level record
represents one of the longest continuous, managed groundwater
observation series in North America. The system features explicit,
legally codified regulatory thresholds (the 660-ft Stage I drought
trigger defined by the Edwards Aquifer Authority, EAA) as well as
critical physical tipping points (the approximately 618-ft head level at
which Comal Springs, the largest freshwater discharge point in the
southwestern United States, approaches total cessation). Furthermore,
continuous historical estimates of annual regional recharge and
groundwater withdrawals are available, constructed independently of the
water-level record.

Because regional flow through this prolific karst system has
historically been modeled using both equivalent porous media (EPM)
approximations and lumped multi-compartment schemes (Scanlon et al.,
2003), evaluating the empirical forecast limits of a single-pool
water-balance mapping addresses a long-standing hydrogeological debate.

This investigation tests whether dynamic stock-flow, autoregressive
error, observational delay, or teleconnection-informed climate modules
systematically reduce forecast errors for annual-mean J-17 water levels
relative to naive persistence and a univariate AR(1) baseline.
Climate-informed recharge variants are evaluated because large-scale
modes of atmospheric variability---specifically the El Niño/Southern
Oscillation (ENSO)---are known to modulate cool-season precipitation
across the southern United States (Ropelewski \& Halpert, 1986).
Pre-season indices, such as the September--November Niño 3.4 sea-surface
temperature anomaly and antecedent Texas climate-division precipitation,
represent operational signals available at the start of a calendar-year
forecast.

A companion study evaluates this identical model-ablation protocol on a
long-lived marine fishery (Northern cod, NAFO Divs. 2J3KL; Abaee, 2026),
testing how structural models behave when dynamic environmental forcing
exhibits low temporal persistence across forecast origins. The
evaluations are kept strictly independent: data are never pooled, and
retention outcomes are not transferred between domains.

\begin{center}\rule{0.5\linewidth}{0.5pt}\end{center}

\subsection{2. Data and System
Specification}\label{data-and-system-specification}

To eliminate ambiguity regarding data inputs, statistical constructs,
predictive models, and regulatory standards, all components are formally
categorized within the San Antonio Pool specification (Table 1).

\subsubsection{Table 1. System specification for the San Antonio Pool of
the Edwards Aquifer and the J-17 index well
(1934--2023).}\label{table-1.-system-specification-for-the-san-antonio-pool-of-the-edwards-aquifer-and-the-j-17-index-well-19342023.}

{\def\LTcaptype{none} % do not increment counter
\begin{longtable}[]{@{}
  >{\raggedright\arraybackslash}p{(\linewidth - 4\tabcolsep) * \real{0.3077}}
  >{\raggedright\arraybackslash}p{(\linewidth - 4\tabcolsep) * \real{0.3077}}
  >{\centering\arraybackslash}p{(\linewidth - 4\tabcolsep) * \real{0.3846}}@{}}
\toprule\noalign{}
\begin{minipage}[b]{\linewidth}\raggedright
Specification Field
\end{minipage} & \begin{minipage}[b]{\linewidth}\raggedright
Description and Archival Source
\end{minipage} & \begin{minipage}[b]{\linewidth}\centering
Variable Type
\end{minipage} \\
\midrule\noalign{}
\endhead
\bottomrule\noalign{}
\endlastfoot
\textbf{Physical System} & Edwards (Balcones Fault Zone) Aquifer, San
Antonio Pool, indexed by monitoring well J-17 & Data (\(\mathrm{D}\)) \\
\textbf{Target Predictand (\(H_t\))} & Calendar-year mean of daily
maximum groundwater elevations (ft above mean sea level, AMSL) & Data
(\(\mathrm{D}\)) \\
\textbf{Observation Point} & TWDB State Well 68-37-203 / EAA Well
AY-68-37-203, Bexar County, Texas & Data (\(\mathrm{D}\)) \\
\textbf{Spatial / Temporal Domain} & San Antonio Pool; continuous
calendar-year series covering 1934--2023 (\(N = 90\text{ years}\)) &
Data (\(\mathrm{D}\)) \\
\textbf{Physical Service Boundary} & Hydraulic head sufficient to
sustain perennial discharge at Comal Springs (\(\approx 618\text{ ft}\);
1956 daily minimum: \(612.51\text{ ft}\)) & Normative / Empirical
(\(\mathrm{N / E}\)) \\
\textbf{Regulatory Boundary} & Statutory Stage I Critical Period
Trigger: 10-day running mean at J-17 \(< 660\text{ ft}\) AMSL &
Normative (\(\mathrm{N}\)) \\
\textbf{Stochastic Disturbance} & High-magnitude unconfined recharge
pulses; conduit-matrix exchanges; unmeasured Uvalde Pool inter-pool
transfer & Model (\(\mathrm{M}\)) \\
\textbf{Groundwater Extraction (\(P_t\))} & Total annual regional well
withdrawals: pre-regulatory agricultural/municipal pumping and post-1996
EAA-permitted withdrawals & Empirical (\(\mathrm{E}\)) \\
\textbf{Ecological Service Record} & USGS Station 08168710 Comal Springs
at New Braunfels, Texas, annual mean discharge (\(\text{cfs}\)) & Data
(\(\mathrm{D}\)) \\
\end{longtable}
}

\emph{Note to Table 1:} Field typing distinguishes direct observations
(\(\mathrm{D}\)), empirical constructs derived from regional balances
(\(\mathrm{E}\)), structural model assumptions (\(\mathrm{M}\)), and
normative regulatory or physical thresholds (\(\mathrm{N}\)).

The target predictand \(H_t\) is the calendar-year arithmetic mean of
daily high water-level elevations recorded at J-17. The series is
derived from continuous automated recordings maintained by the Texas
Water Development Board (TWDB) and the EAA. The historical record
includes the pre-continuous Beverly Lodges measurement composite
spanning 1934--1962, which is incorporated as officially published.

To maintain sampling consistency across the 90-year panel, years with
fewer than 240 valid daily observations are excluded. This threshold
never binds on the analyzed series: the lowest observation counts occur
in 1939 (\(n = 242\)) and 1935 (\(n = 258\)), both of which are retained
without synthetic daily interpolation. The terminal year (2023) reflects
provisional TWDB status. Evaluating the mean of daily highs provides a
stable annual index; choosing a daily-mean convention would shift the
baseline slightly downward, which is relevant primarily when
interpreting trajectories near the 660-ft regulatory trigger and the
618-ft physical springflow cessation limit.

Historical recharge (\(R_t\)) represents total annual estimated recharge
for the San Antonio Pool (\(10^3\text{ acre-ft}\cdot\text{yr}^{-1}\)),
calculated by the U.S. Geological Survey (USGS) using the stream-loss
water-budget methodology (Puente, 1978; Umphres \& Choi, 2025). This
estimate is constructed from basin precipitation, runoff modeling, and
upstream/downstream streamflow gauging on nine major catchment streams
crossing the Edwards Aquifer recharge zone, entirely independent of
water-level observations at J-17.

Historical groundwater extraction (\(P_t\)) is total annual well
discharge (\(10^3\text{ acre-ft}\cdot\text{yr}^{-1}\)) compiled in EAA
historical summary tables (Edwards Aquifer Authority, 2024/25). Annual
discharge from Comal Springs (\(Q_t\), in cubic feet per second,
\(\text{cfs}\)) is obtained from USGS Station 08168710; it is held as a
secondary service series and deliberately excluded as an input to the
predictive head equations.

Antecedent climate variables comprise: * The September--November average
sea-surface temperature anomaly in the Niño 3.4 region
(\(5^\circ\text{N}–5^\circ\text{S}, 170^\circ\text{W}–120^\circ\text{W}\))
from the HadISST dataset, calculated relative to a 1991--2020
climatology. * Total calendar-year precipitation across Texas Climate
Divisions 06 (Edwards Plateau) and 07 (South Central) derived from the
NOAA NCEI nClimDiv database.

\begin{center}\rule{0.5\linewidth}{0.5pt}\end{center}

\subsection{3. Forecast Models}\label{forecast-models}

The evaluated models form a forward-ordered structural ladder (Table 2).
Each structural rung adds a specific hydrogeological mechanism to the
preceding configuration, moving from naive statistical rules to dynamic
stock-flow equations, residual autoregressions, and climate-informed
teleconnections.

\subsubsection{Definition 3.1 (Lumped One-Pool Water-Balance
Map)}\label{definition-3.1-lumped-one-pool-water-balance-map}

Let \(H_t\) denote the annual-mean water level at index well J-17 at the
end of calendar year \(t\). The discrete-time one-pool water-balance
mapping is defined as:
\[H_{t+1} = \Big[ H_t + \alpha + \beta \tilde{R}_{t+1} + \gamma \tilde{P}_{t+1} + \delta H_t \Big]_{\mathrm{clip}}\]
where
\([\,\cdot\,]_{\mathrm{clip}} = \min(\max(\,\cdot\,, 610), 710)\text{ ft AMSL}\),
and \(\tilde{R}_{t+1}\) and \(\tilde{P}_{t+1}\) denote the proxy flux
values substituted for the unknown future recharge and extraction
occurring across year \(t+1\).

The model coefficients are estimated via ordinary least squares (OLS)
over the available training window: * \(\alpha\) represents an intercept
term; * \(\beta\) represents head response per unit of regional
recharge; * \(\gamma\) represents head drawdown per unit of well
extraction; * \(\delta\) represents a linear regional drainage
coefficient governing self-decay and baseflow discharge.

\subsubsection{Table 2. Forward-ordered model ladder and structural
specifications.}\label{table-2.-forward-ordered-model-ladder-and-structural-specifications.}

{\def\LTcaptype{none} % do not increment counter
\begin{longtable}[]{@{}
  >{\raggedright\arraybackslash}p{(\linewidth - 6\tabcolsep) * \real{0.2500}}
  >{\raggedright\arraybackslash}p{(\linewidth - 6\tabcolsep) * \real{0.2500}}
  >{\raggedright\arraybackslash}p{(\linewidth - 6\tabcolsep) * \real{0.2500}}
  >{\raggedright\arraybackslash}p{(\linewidth - 6\tabcolsep) * \real{0.2500}}@{}}
\toprule\noalign{}
\begin{minipage}[b]{\linewidth}\raggedright
Model ID
\end{minipage} & \begin{minipage}[b]{\linewidth}\raggedright
Structural Class
\end{minipage} & \begin{minipage}[b]{\linewidth}\raggedright
Projected Flux Assumptions (\(\tilde{R}_{t+h}, \tilde{P}_{t+h}\))
\end{minipage} & \begin{minipage}[b]{\linewidth}\raggedright
Operational Role in Evaluation
\end{minipage} \\
\midrule\noalign{}
\endhead
\bottomrule\noalign{}
\endlastfoot
\textbf{persist} & Naive Baseline & None & \(\hat{H}_{t+h} = H_t\) \\
\textbf{mean} & Climatological Baseline & None &
\(\hat{H}_{t+h} = \bar{H}_{\mathrm{train}}\) \\
\textbf{M1} & Autonomous AR(1) & None &
\(\hat{H}_{t+1} = a + \varphi H_t\) \\
\textbf{M2} & Causal Stock-Flow & Persisted prior-year fluxes:
\(\tilde{R}_{t+h} = R_t\), \(\tilde{P}_{t+h} = P_t\) & Candidate
structural model \\
\textbf{M2m} & Climatological Stock-Flow & Training-window mean fluxes:
\(\tilde{R}_{t+h} = \bar{R}\), \(\tilde{P}_{t+h} = \bar{P}\) &
Structurally affine AR(1) \\
\textbf{M3} & Residual Autoregressive & Persisted fluxes; residual
correction: \(\hat{e}_{t+k} = \hat{\theta}^k e_t\) & AR(1) tracking on
head errors \\
\textbf{M4} & Lagged Initialization & Persisted fluxes; initiated from
stale state: \(\hat{H}_{t+1 \mid t} = f(H_{t-1})\) & Theoretical
information-lag control \\
\textbf{M2\_oracle} & Diagnostic Ceiling & Realized future fluxes:
\(\tilde{R}_{t+h} = R_{t+h}\), \(\tilde{P}_{t+h} = P_{t+h}\) & Benchmark
ceiling (excluded from retention) \\
\end{longtable}
}

\emph{Role note.} The delay-module M4 is retained as a symmetry control
with the companion fisheries evaluation (Abaee 2026b), not as a
monitoring constraint of this system; in the ladder mapping, M4 inherits
the flux regime of M2 (``as M2'') by construction.

Covariates enter each structural rung on a
last-observation-carried-forward convention: the most recent observation
at or before the forecast origin is used, and no covariate value dated
after the origin enters any fitted transition.

\subsubsection{Remark 3.1 (Class Reduction of Model
M2m)}\label{remark-3.1-class-reduction-of-model-m2m}

Under stationary, climatological flux assumptions
(\(\tilde{R}_{t+1} = \bar{R}\) and \(\tilde{P}_{t+1} = \bar{P}\)), the
dynamic one-pool balance simplifies algebraically to:
\[H_{t+1} = (1 + \delta) H_t + (\alpha + \beta \bar{R} + \gamma \bar{P})\]
Defining \(\varphi = 1 + \delta\) and
\(a = \alpha + \beta \bar{R} + \gamma \bar{P}\), Model M2m reduces
functionally to the identical mathematical class as the autonomous
autoregressive model M1 (\(H_{t+1} = \varphi H_t + a\)). However, M2m
uses a distinct estimator: rather than fitting a univariate regression
directly on historical head levels, it pins its intercept using
in-sample flux averages and least-squares response coefficients.

As detailed in Section 4.1, M2m is formally declined within our protocol
based on this structural equivalence, rather than being certified as a
distinct causal stock-flow innovation.

In auxiliary climate-informed variants (Section 5.4), persisted recharge
\(\tilde{R}_{t+1} = R_t\) is replaced by an out-of-sample statistical
prediction \(\hat{R}_{t+1}\) generated from information available at the
December 31 forecast origin: * \textbf{M2\_Rar:} First-order
autoregression fit to the historical recharge series; *
\textbf{M2\_Renso:} Linear regression forced by antecedent autumn Niño
3.4 SST anomalies; * \textbf{M2\_Rprecip:} Linear regression forced by
contemporaneous calendar-year Texas Division 06 and 07 precipitation; *
\textbf{M2\_combo:} Multivariate regression combining antecedent
recharge, Niño 3.4 anomalies, and division precipitation.

\begin{center}\rule{0.5\linewidth}{0.5pt}\end{center}

\subsection{4. Evaluation Design}\label{evaluation-design}

To eliminate retrospective selection bias, the quantitative evaluation
protocol and retention gates were formally locked prior to generating
predictive scores.

\subsubsection{Definition 4.1 (Primary and Secondary Evaluation
Metrics)}\label{definition-4.1-primary-and-secondary-evaluation-metrics}

The primary evaluation metric is the out-of-sample Root Mean Squared
Error (RMSE) of annual-mean J-17 elevations, expressed in feet:
\[\mathrm{RMSE}_h = \sqrt{\frac{1}{n} \sum_{i=1}^n \left( \hat{H}_{t_i+h \mid t_i} - H_{t_i+h} \right)^2}\]
Secondary diagnostic metrics comprise: * \textbf{Mean Absolute Error
(MAE):} Expressed in feet. * \textbf{Brier Score on Regulatory
Compliance:} The mean squared error of the deterministic threshold
indicator \(\mathbf{1}\{\hat{H}_{t+h} < 660\text{ ft}\}\), evaluated
over the post-2007 regulatory period.

\subsubsection{Definition 4.2 (The Frozen Retention
Rule)}\label{definition-4.2-the-frozen-retention-rule}

Let \(\mathrm{RMSE}(M)\) denote the primary rolling-origin RMSE for
candidate model \(M\) at forecast horizon \(h=1\). A candidate causal
module \(M\) from the structural ladder is retained if and only if it
satisfies two conditions: 1. \textbf{Hypothesis 1 (Baseline Superiority,
H1):} \(\mathrm{RMSE}(M) < \mathrm{RMSE}(\text{persist})\). 2.
\textbf{Hypothesis 2 (Parsimony Progression, H2):}
\(\mathrm{RMSE}(M) < \mathrm{RMSE}(M_{\text{next-simpler causal}})\).

Retention is determined strictly at the one-year-ahead horizon (\(h=1\))
based on the point-RMSE ranking without an initial tie tolerance.
Diagnostic oracles (M2\_oracle) and the auxiliary Comal Springs
discharge series are barred from the retention competition.

\subsubsection{4.1 Protocol Record and Pre-Declared
Deviations}\label{protocol-record-and-pre-declared-deviations}

The primary evaluation protocol was frozen and archived on 2026-08-25
prior to executing model scoring. Three procedural deviations from the
baseline text of Definition 4.2 are documented: 1. \textbf{The M2m Class
Exclusion:} Model M2m satisfies both (H1) and (H2) at \(h=1\) on
empirical point scores. However, because it collapses algebraically to
an affine AR(1) process (Remark 3.1), an explicit protocol clause
excludes it from retention as an independent causal stock-flow model. 2.
\textbf{The Climate-Rung Comparator:} The protocol specification for
Pass 2 evaluates whether climate-informed recharge predictions improve
head forecasts relative to persistence and relative to the retained
model M1. Using the declined M2m model as the causal comparator for the
climate rungs would be circular; we therefore evaluate the climate
models against both M1 and persistence, retaining M2m as a nested
reference benchmark. 3. \textbf{Directional Sign-Hit Metric:} A sign-hit
rate metric for annual increments (\(\Delta H\)) included in initial
exploratory outlines was struck prior to scoring because persistence
generates \(\Delta H = 0\) (rendering directional comparisons
mathematically undefined).

Fixed evaluation windows are designed around major historical climate
and regulatory transitions: * \textbf{Drought-of-Record Drawdown:}
Training 1934--1950 (\(n=16\) transitions), testing 1951--1956 (\(n=6\)
years). * \textbf{Drought-of-Record Recovery:} Training 1934--1956
(\(n=22\) transitions), testing 1957--1961 (\(n=5\) years). *
\textbf{Pre-Permit Wet Era:} Training 1980--1990 (\(n=10\) transitions),
testing 1991--1995 (\(n=5\) years). * \textbf{Modern Critical-Period
Management:} Training 1997--2014 (\(n=17\) transitions), testing
2015--2023 (\(n=9\) years).

The rolling-origin cross-validation uses an expanding training window
with a minimum threshold of 15 years (\(n=14\) transitions). Scoring
across the 1934--2023 domain yields \(n = 75\) origins for the
one-year-ahead horizon (\(h=1\), origins 1948--2022) and \(n = 71\)
origins for the five-year-ahead horizon (\(h=5\), origins 1948--2018).

\begin{center}\rule{0.5\linewidth}{0.5pt}\end{center}

\subsection{5. Results}\label{results}

\subsubsection{5.1 Historical Water-Level Dynamics at
J-17}\label{historical-water-level-dynamics-at-j-17}

!{[}Figure 1. Annual-mean groundwater elevations at index well J-17
across the 1934--2023 observation period. The 1956 drought-of-record
minimum reached an annual mean of \(623.15\text{ ft}\) (daily minimum
\(612.51\text{ ft}\)); the 1992 pluvial maximum reached an annual mean
of \(691.96\text{ ft}\) (daily maximum \(703.31\text{ ft}\)). Across the
90-year panel, the annual mean falls below the 660-ft regulatory trigger
in 31 years.* --- \#\#\# 5.2 Out-of-Sample Performance on Fixed
Historical Windows Multi-step forward trajectories across the four
diagnostic windows highlight fundamental behavioral differences between
the models (Table 3 and Figure 2). \#\#\# Table 3. Out-of-sample
forecast RMSE (ft) across fixed historical evaluation windows.
\emph{Note: Bold font denotes the best performance among candidate
predictive models within each window. The diagnostic oracle (M2\_oracle)
is shown for reference.} \textbar{} Historical Evaluation Window
\textbar{} persist \textbar{} mean \textbar{} M1 \textbar{} M2
\textbar{} M2m \textbar{} M3 \textbar{} M4 \textbar{} M2\_oracle
\textbar{} \textbar{} :--- \textbar{} :---: \textbar{} :---: \textbar{}
:---: \textbar{} :---: \textbar{} :---: \textbar{} :---: \textbar{}
:---: \textbar{} :---: \textbar{} \textbar{} \textbf{Drawdown
(1951--1956)} \textbar{} 23.75 \textbar{} 35.19 \textbar{} 30.94
\textbar{} \textbf{18.11} \textbar{} 27.44 \textbar{} 18.12 \textbar{}
18.23 \textbar{} 19.69 \textbar{} \textbar{} \textbf{Recovery
(1957--1961)} \textbar{} 43.62 \textbar{} 14.07 \textbar{} 56.24
\textbar{} 55.32 \textbar{} 37.74 \textbar{} 55.28 \textbar{} 55.12
\textbar{} \textbf{12.26} \textbar{} \textbar{} \textbf{Pre-Permit Wet
(1991--1995)} \textbar{} 30.13 \textbar{} 18.24 \textbar{} 20.02
\textbar{} 16.67 \textbar{} 23.47 \textbar{} 16.41 \textbar{}
\textbf{15.26} \textbar{} \textbf{7.18} \textbar{} \textbar{}
\textbf{Critical-Period (2015--2023)} \textbar{} 27.41 \textbar{} 14.77
\textbar{} 15.62 \textbar{} 23.37 \textbar{} 14.79 \textbar{} 22.84
\textbar{} 22.17 \textbar{} \textbf{8.70} \textbar{}

\pandocbounded{\includegraphics[keepaspectratio,alt={Figure 2. Forecast trajectories across historical drought and recovery phases. Persisting drought recharge causes M2 to collapse into the lower clipping boundary (610\textbackslash text\{ ft\}) during the 1957--1961 recovery.}]{figs/fig1_series.png}}
On the 1951--1956 drought-of-record drawdown window, Model M2 achieves
the lowest predictive error (\(18.11\text{ ft}\)). However, this
apparent predictive success is an artifact of continuing drought
persistence rather than hydrogeological realism. Fitting the
four-parameter balance on the brief 1934--1950 baseline (\(n=16\)
transitions) yields a positive extraction coefficient
(\(\hat{\gamma} = +0.021\text{ ft}\cdot(10^3\text{ acre-ft})^{-1}\)),
which violates mass-conservation principles. As the historical drought
developed, regional agricultural and municipal pumping surged in
response to diminished rainfall. Ordinary least squares aliases this
human behavioral response directly into the state equation---an example
of simultaneity bias compounded by small-sample limitations. Because the
terminal training year (1950) was already characterized by depressed
recharge, persisting that low recharge value through the test window
happens to project a continuing decline, outperforming even the
diagnostic oracle (\(19.69\text{ ft}\)). Conversely, Model M2 fails
completely during the 1957--1961 post-drought recovery. The terminal
drought year (1956) saw regional recharge collapse to
\(43.7 \times 10^3\text{ acre-ft}\), followed by a major pluvial shift
in 1957 (\(1,143 \times 10^3\text{ acre-ft}\)). Persisting the dry 1956
flux causes M2 to project an escalating collapse: projected head drops
from \(617.1\text{ ft}\) to \(611.6\text{ ft}\), \(606.8\text{ ft}\),
and ultimately saturates against the artificial lower boundary clip at
\(610.0\text{ ft}\) from 1959 onward, generating an RMSE of
\(55.32\text{ ft}\). Naive persistence performs nearly as poorly
(\(43.62\text{ ft}\)) by freezing the depleted 1956 head level
(\(623.15\text{ ft}\)). In contrast, the diagnostic oracle (M2\_oracle),
which receives the realized 1957 flood recharge, tracks the water-level
surge (\(12.26\text{ ft}\)). During rapid pluvial recoveries, long-term
training climatology provides the best non-oracle forecast
(\(14.07\text{ ft}\)), demonstrating the strong mean-reverting behavior
of this karst aquifer. --- \#\#\# 5.3 Rolling-Origin Forecast Evaluation
Across 75 expanding rolling origins, naive persistence achieves an RMSE
of \(13.23\text{ ft}\) at \(h=1\) and \(21.11\text{ ft}\) at \(h=5\)
(Tables 4 and 5, Figure 3). \#\#\# Table 4. Rolling-origin forecast
performance across the 1934--2023 panel (ft).{]}(figs/fig2\_windows.png)

Note: Evaluated across identical origin sets (\(n=75\) for \(h=1\);
\(n=71\) for \(h=5\)). Bold font indicates the best non-oracle
forecast.* \textbar{} Model ID \textbar{} One-Year RMSE (\(h=1\))
\textbar{} One-Year MAE (\(h=1\)) \textbar{} Five-Year RMSE (\(h=5\))
\textbar{} \textbar{} :--- \textbar{} :---: \textbar{} :---: \textbar{}
:---: \textbar{} \textbar{} \textbf{persist} \textbar{} 13.23 \textbar{}
10.73 \textbar{} 21.11 \textbar{} \textbar{} \textbf{mean} \textbar{}
16.17 \textbar{} 13.17 \textbar{} \textbf{16.80} \textbar{} \textbar{}
\textbf{M1} \textbar{} 12.84 \textbar{} 10.72 \textbar{} 21.25
\textbar{} \textbar{} \textbf{M2} \textbar{} 14.70 \textbar{} 11.45
\textbar{} 33.49 \textbar{} \textbar{} \textbf{M2m} \textbar{}
\textbf{12.28} \textbar{} \textbf{10.22} \textbar{} 17.44 \textbar{}
\textbar{} \textbf{M3} \textbar{} 14.46 \textbar{} 11.12 \textbar{}
33.46 \textbar{} \textbar{} \textbf{M4} \textbar{} 14.30 \textbar{}
11.17 \textbar{} 33.39 \textbar{} \textbar{} \textbf{M2\_oracle}
\textbar{} 7.55 \textbar{} 5.79 \textbar{} 10.87 \textbar{}

\subsubsection{\texorpdfstring{Table 5. Formal retention evaluation on
primary rolling-origin RMSE at
\(h=1\).}{Table 5. Formal retention evaluation on primary rolling-origin RMSE at h=1.}}\label{table-5.-formal-retention-evaluation-on-primary-rolling-origin-rmse-at-h1.}

{\def\LTcaptype{none} % do not increment counter
\begin{longtable}[]{@{}
  >{\raggedright\arraybackslash}p{(\linewidth - 8\tabcolsep) * \real{0.1739}}
  >{\centering\arraybackslash}p{(\linewidth - 8\tabcolsep) * \real{0.2174}}
  >{\centering\arraybackslash}p{(\linewidth - 8\tabcolsep) * \real{0.2174}}
  >{\centering\arraybackslash}p{(\linewidth - 8\tabcolsep) * \real{0.2174}}
  >{\raggedright\arraybackslash}p{(\linewidth - 8\tabcolsep) * \real{0.1739}}@{}}
\toprule\noalign{}
\begin{minipage}[b]{\linewidth}\raggedright
Model ID
\end{minipage} & \begin{minipage}[b]{\linewidth}\centering
Margin vs.~Persistence (ft)
\end{minipage} & \begin{minipage}[b]{\linewidth}\centering
Margin vs.~Simpler Comparator (ft)
\end{minipage} & \begin{minipage}[b]{\linewidth}\centering
Distinct Structural Class?
\end{minipage} & \begin{minipage}[b]{\linewidth}\raggedright
Protocol Retention Verdict
\end{minipage} \\
\midrule\noalign{}
\endhead
\bottomrule\noalign{}
\endlastfoot
\textbf{M1} & \(-0.39\) (\(12.84 < 13.23\)) & None (Base model) &
Output-only linear AR(1) & \textbf{Retained} (Point rule; margin within
noise) \\
\textbf{M2} & \(+1.47\) (\(14.70 > 13.23\)) & \(+1.86\) vs.~M1 & Dynamic
stock-flow & \textbf{Rejected} (Fails H1) \\
\textbf{M2m} & \(-0.95\) (\(12.28 < 13.23\)) & \(-0.56\) vs.~M1 & No
(Collapses to affine AR(1)) & \textbf{Declined} (Excluded by protocol
class clause) \\
\textbf{M3} & \(+1.23\) (\(14.46 > 13.23\)) & \(-0.24\) vs.~M2 &
Residual autoregression & \textbf{Rejected} (Fails H1) \\
\textbf{M4} & \(+1.07\) (\(14.30 > 13.23\)) & \(-0.16\) vs.~M3 & Lagged
state update & \textbf{Rejected} (Fails H1) \\
\textbf{M2\_oracle} & \(-5.68\) (\(7.55 < 13.23\)) & None & Requires
realized future fluxes & \textbf{Excluded} (Diagnostic benchmark
ceiling) \\
\end{longtable}
}

!{[}Figure 3. Out-of-sample RMSE for candidate models at one-year and
five-year forecast horizons. At \(h=1\), M2m and M1 edge out
persistence; at \(h=5\), climatological mean baselines decisively
outperform all dynamic models.* Applying the frozen retention rule
(Definition 4.2) at \(h=1\): * \textbf{Model M1} yields an RMSE of
\(12.84\text{ ft}\), beating persistence by \(0.39\text{ ft}\). Because
M1 represents the initial autonomous baseline, Condition (H2) is
vacuous, and M1 is retained under the point rule. * \textbf{Model M2}
yields an RMSE of \(14.70\text{ ft}\), failing Condition (H1) by
\(+1.47\text{ ft}\) relative to persistence and trailing M1 by
\(+1.86\text{ ft}\). * \textbf{Models M3 and M4} achieve minor
reductions relative to M2 (\(14.46\text{ ft}\) and \(14.30\text{ ft}\)),
but remain inferior to naive persistence, failing Condition (H1). *
\textbf{Model M2m} achieves the lowest one-step error
(\(12.28\text{ ft}\)), outperforming persistence by \(0.95\text{ ft}\)
and beating M1 by \(0.56\text{ ft}\). However, as established in Remark
3.1, M2m collapses to an affine AR(1) state equation and is declined
under the pre-declared protocol clause. The formal verdict that ``no
stock-flow module is retained'' is partly a consequence of this class
exclusion: the water-balance formulation forced by climatological mean
fluxes is the best one-step forecaster evaluated. The primary
operational failure lies in persisting last year's recharge: doing so
inflates the one-year error of M2 by \(+2.42\text{ ft}\) relative to
M2m. In multi-step projections (\(h=5\)), this dynamic error compounds
severely: M2, M3, and M4 diverge, posting errors of approximately
\(33.5\text{ ft}\), while training-mean climatology achieves the lowest
RMSE (\(16.80\text{ ft}\)), outperforming persistence
(\(21.11\text{ ft}\)) by \(4.31\text{ ft}\). \#\#\#\# Secondary
Regulatory Performance (Post-2007 Era) Over the post-2007 regulatory
period (\(n=16\) origins at \(h=1\)), the Brier misclassification score
for the 660-ft Stage I drought trigger is \(0.31\) for persistence,
\(0.25\) for M1, and \(0.19\) for M2\_oracle. At this one-year horizon,
the model ranking is preserved: persistence posts an RMSE of
\(13.09\text{ ft}\), M1 achieves \(12.16\text{ ft}\), M2 yields
\(13.31\text{ ft}\), and the oracle records \(8.03\text{ ft}\). At
\(h=5\) (\(n=12\) origins), the relative ordering between M1 and
persistence reverses on this subset (\(17.16\text{ ft}\)
vs.~\(25.10\text{ ft}\)). These calculations evaluate an annual-mean
proxy against a statutory 10-day trigger on a small subsample,
confirming that the primary point-RMSE findings are not contradicted by
threshold metrics. --- \#\#\# 5.3.1 Statistical Uncertainty Across Model
Retention Margins To assess whether the small performance margins
represent statistically significant improvements, we compute
Diebold--Mariano loss-differential tests (Diebold \& Mariano, 1995)
using Newey--West heteroskedasticity- and autocorrelation-consistent
(HAC) standard errors (lag \(h-1\)) paired with moving-block bootstrap
confidence intervals (Künsch, 1989; 10,000 replications, block length
\(L = 8\); Table 6). \#\#\# Table 6. Statistical evaluation of primary
model retention margins. \emph{Note: Evaluated across rolling origins
(\(n=75\) for \(h=1\); \(n=71\) for \(h=5\)). Block bootstrap uses
10,000 replications with block length \(L=8\). Bold text denotes margins
whose 95\% confidence intervals exclude zero.} \textbar{} Candidate
Comparison Model \textbar{} Horizon (\(h\)) \textbar{} RMSE Margin (ft)
\textbar{} Diebold--Mariano (\(z\)) \textbar{} DM \(p\)-value \textbar{}
95\% Moving-Block Bootstrap CI (ft) \textbar{} Separates from Zero?
\textbar{} \textbar{} :--- \textbar{} :---: \textbar{} :---: \textbar{}
:---: \textbar{} :---: \textbar{} :---: \textbar{} :---: \textbar{}
\textbar{} \textbf{M1 \(-\) persist} \textbar{} 1 \textbar{} \(-0.39\)
\textbar{} \(-0.85\) \textbar{} 0.40 \textbar{} {[}\(-1.51, +0.71\){]}
\textbar{} No \textbar{} \textbar{} \textbf{M2m \(-\) persist}
\textbar{} 1 \textbar{} \(-0.95\) \textbar{} \(-3.07\) \textbar{} 0.003
\textbar{} {[}\textbf{\(-1.45, -0.68\)}{]} \textbar{} \textbf{Yes}
\textbar{} \textbar{} \textbf{M2m \(-\) M1} \textbar{} 1 \textbar{}
\(-0.56\) \textbar{} \$-1. The M2 minus persistence margin of
\(+1.47\text{ ft}\) remains within noise under the post-freeze layer (DM
statistic \(+1.27\), \(p = 0.21\), 95\% block-bootstrap CI
\([-0.02, +3.40]\text{ ft}\)).61\$ \textbar{} 0.11 \textbar{}
{[}\(-1.53, +0.13\){]} \textbar{} No \textbar{} \textbar{} \textbf{M2
\(-\) persist} \textbar{} 1 \textbar{} \(+1.47\) \textbar{} \(+1.27\)
\textbar{} 0.21 \textbar{} {[}\(-0.02, +3.40\){]} \textbar{} No
(Borderline) \textbar{} \textbar{} \textbf{mean \(-\) persist}
\textbar{} 5 \textbar{} \(-4.30\) \textbar{} \(-1.65\) \textbar{} 0.10
\textbar{} {[}\textbf{\(-9.72, -2.09\)}{]} \textbar{} \textbf{Yes}
\textbar{} \textbar{} \textbf{M1 \(-\) persist} \textbar{} 5 \textbar{}
\(+0.15\) \textbar{} \(+0.06\) \textbar{} 0.95 \textbar{}
{[}\(-5.57, +4.94\){]} \textbar{} No \textbar{} \textbar{}
\textbf{M2\_combo \(-\) M1} \textbar{} 1 \textbar{} \(-0.13\) \textbar{}
\(-0.21\) \textbar{} 0.83 \textbar{} {[}\(-1.37, +1.31\){]} \textbar{}
No \textbar{} \textbar{} \textbf{M2\_combo \(-\) M2m} \textbar{} 1
\textbar{} \(+0.43\) \textbar{} \(+0.74\) \textbar{} 0.46 \textbar{}
{[}\(-0.56, +1.57\){]} \textbar{} No \textbar{} The uncertainty analysis
provides three key insights: 1. \textbf{The M1 Retention Margin is
Within Noise:} The \(0.39\text{ ft}\) error reduction achieved by M1
over persistence yields a Diebold--Mariano statistic of \(z = -0.85\)
(\(p = 0.40\)), and its 95\% bootstrap confidence interval
({[}\(-1.51, +0.71\){]} ft) spans zero. The mean absolute errors of the
two rules are essentially identical (\(10.72\text{ ft}\)
vs.~\(10.73\text{ ft}\)). Retaining M1 under the point rule simply
captures a slight empirical mean reversion in head levels, rather than a
decisive predictive advantage. 2. \textbf{M2m is the Sole Statistically
Separable One-Step Model:} The \(0.95\text{ ft}\) advantage of M2m over
persistence is highly statistically significant (\(z = -3.07\),
\(p = 0.003\); 95\% bootstrap CI: {[}\(-1.45, -0.68\){]} ft). However,
the \(0.56\text{ ft}\) increment between M2m and M1 is not statistically
distinguishable from zero (\(p = 0.11\), CI: {[}\(-1.53, +0.13\){]} ft).
The estimator-level benefit of incorporating mean fluxes is observable
in point estimates, but remains within sampling uncertainty at this
sample size. 3. \textbf{Climatology Decisively Beats Persistence at Long
Horizons:} At \(h=5\), the training mean's \(-4.30\text{ ft}\) advantage
over persistence yields a bootstrap interval that strictly excludes zero
({[}\(-9.72, -2.09\){]} ft). While the multi-step Diebold--Mariano test
exhibits reduced power due to overlapping prediction horizons
(\(p = 0.10\)), the bootstrap distribution confirms that climatological
mean reversion reliably outperforms persistence over multi-year spans.
An independent computational audit of these calculations (using an
unweighted truncation scheme, population variance scaling, block length
\(L=h\), independent seeds, and alternative HAC conventions) confirmed
every empirical finding: M1 vs.~persistence yielded \(z = -0.86\)
(bootstrap CI: {[}\(-1.22, +0.56\){]} ft; \(p = 0.38\)), M2m
vs.~persistence yielded \(p = 0.001\), and the five-year climatology
advantage remained significant (\(p = 0.035\)). --- \#\#\# 5.4
Climate-Informed Recharge Predictions We next evaluate whether coupling
the one-pool balance to pre-season climate teleconnections can overcome
the timing errors inherent in persisted recharge (Table 7). \#\#\# Table
7. Rolling-origin forecast skill for climate-informed recharge modules.
\emph{Note: Evaluated across rolling origins (\(n=75\) for \(h=1\);
\(n=71\) for \(h=5\)). The precipitation oracle uses realized year
\(t+h\) climate-division precipitation and is excluded from retention.}
\textbar{} Model ID \textbar{} Projected Recharge Input Formulation
(\(\tilde{R}_{t+1}\)) \textbar{} Head RMSE \(h=1\) (ft) \textbar{} Head
Margin vs.~M1 (ft) \textbar{} Head RMSE \(h=5\) (ft) \textbar{} Recharge
Target RMSE (\(10^3\text{ acre-ft}\)) \textbar{} \textbar{} :---
\textbar{} :--- \textbar{} :---: \textbar{} :---: \textbar{} :---:
\textbar{} :---: \textbar{} \textbar{} \textbf{persist} \textbar{}
Persisted head / Persisted recharge (\(R_t\)) \textbar{} 13.23
\textbar{} \(+0.39\) \textbar{} \textbf{21.11} \textbar{} 702 \textbar{}
\textbar{} \textbf{M1} \textbar{} None (Autonomous linear AR(1))
\textbar{} 12.84 \textbar{} --- \textbar{} 21.25 \textbar{} ---
\textbar{} \textbar{} \textbf{M2\_Rar} \textbar{} Univariate AR(1) on
historical recharge \textbar{} 13.25 \textbar{} \(+0.41\) \textbar{}
25.38 \textbar{} 561 \textbar{} \textbar{} \textbf{M2\_Renso} \textbar{}
Antecedent Sep--Nov Niño 3.4 SST anomaly \textbar{} 12.82 \textbar{}
\(-0.02\) \textbar{} 24.42 \textbar{} 528 \textbar{} \textbar{}
\textbf{M2\_Rprecip} \textbar{} Prior-year Texas Div 06/07 precipitation
\textbar{} 12.80 \textbar{} \(-0.04\) \textbar{} 25.38 \textbar{} 545
\textbar{} \textbar{} \textbf{M2\_combo} \textbar{} Multivariate: AR(1)
+ Niño 3.4 + Precip \textbar{} 12.71 \textbar{} \(-0.13\) \textbar{}
26.88 \textbar{} 538 \textbar{} \textbar{} \textbf{M2m} \textbar{}
Climatological mean recharge (\(\bar{R}_{\mathrm{train}}\)) \textbar{}
\textbf{12.28} \textbar{} \textbf{\(-0.56\)} \textbar{} 17.44 \textbar{}
556 \textbar{} \textbar{} \textbf{rain\_oracle} \textbar{} Realized year
\(t+1\) Div 06/07 precipitation \textbar{} 10.56 \textbar{} \(-2.28\)
\textbar{} 16.91 \textbar{} \textbf{354} \textbar{} Contemporaneous
precipitation is strongly correlated with annual recharge
(\(\operatorname{corr}(R_t, \text{Precip}_t) = 0.78\)). Consequently, if
an analyst had perfect foreknowledge of annual rainfall (the
rain\_oracle), head RMSE would drop to \(10.56\text{ ft}\). However,
this still trails the full diagnostic oracle M2\_oracle
(\(7.55\text{ ft}\)) because a linear rainfall-recharge relationship
fails to capture the extreme non-linear runoff generation observed
during major pluvials such as 1957. Antecedent climate variables known
at the forecast origin (December 31) demonstrate modest predictive skill
for next-year recharge relative to a persistence baseline: M2\_Renso,
M2\_Rprecip, and M2\_combo achieve recharge RMSEs of
\(528\text{--}545 \times 10^3\text{ acre-ft}\), outperforming persisted
recharge (\(702 \times 10^3\text{ acre-ft}\)) and training-mean
climatology (\(556 \times 10^3\text{ acre-ft}\)). However, this modest
improvement in recharge estimation fails to translate into meaningful
predictive skill for index-well head: * The climate modules yield
minimal head RMSE reductions over M1 (\(0.02\text{ ft}\) for M2\_Renso,
\(0.04\text{ ft}\) for M2\_Rprecip, and \(0.13\text{ ft}\) for
M2\_combo). As demonstrated in Table 6, these margins are within
sampling noise (for M2\_combo vs.~M1, 95\% bootstrap CI:
{[}\(-1.37, +1.31\){]} ft). * When evaluated against the nested
climatological baseline M2m (\(12.28\text{ ft}\)), the climate variants
perform worse by \(+0.43\text{ to } +0.54\text{ ft}\). * In multi-step
projections (\(h=5\)), freezing the one-step climate recharge forecast
causes head RMSE to inflate to \(24.42\text{--}26.88\text{ ft}\),
falling far behind both persistence (\(21.11\text{ ft}\)) and
climatology (\(16.80\text{ ft}\)). * An autoregressive model on recharge
itself (M2\_Rar) degrades head forecasts (\(13.25\text{ ft}\) at
\(h=1\)), confirming that fitting an AR(1) process to near-white
recharge noise introduces spurious parameter variance. On fixed
historical windows, antecedent climate signals fail to anticipate abrupt
hydrologic transitions. For example, September--November 1956 is La Niña
(\(-0.92\)) and does not announce
\(R_{1957} = 1143 \times 10^3\text{ acre-ft}\): the teleconnection
carried no forecast value for the record pluvial that ended the drought.
Consequently, all causal climate-informed models generated large errors
across the recovery window (RMSEs of \(48.8\text{--}55.0\text{ ft}\)),
trailing naive persistence (\(43.62\text{ ft}\)) and the rain oracle
(\(33.72\text{ ft}\)). ---
\pandocbounded{\includegraphics[keepaspectratio,alt={Figure 4. Rolling RMSE on J-17 for climate-informed recharge. The precipitation oracle uses year t+h precipitation and cannot be retained.}]{figs/fig4_pass2.png}}

\subsubsection{5.5 Hydrological Redundancy of the Comal Springs Service
Series}\label{hydrological-redundancy-of-the-comal-springs-service-series}

\pandocbounded{\includegraphics[keepaspectratio,alt={Figure 5. Contemporaneous annual Comal Springs discharge (Q\_t) plotted against J-17 water levels (H\_t) over 1934--2023 (\textbackslash operatorname\{corr\}(H\_t, Q\_t) = 0.986). The linear rating curve fit to 1934--1950 data (Q = -2876 + 4.77H) incorrectly predicts positive discharge (97\textbackslash text\{ cfs\}) during the 1956 drought of record.}]{figs/fig3_rmse.png}}
A linear rating curve fit to 1934--1950 data: \[Q_t = -2876 + 4.77 H_t\]
accurately characterizes ordinary operating conditions, reflecting the
high contemporaneous correlation between J-17 water levels and Comal
Springs discharge (\(r = 0.986\); Figure 4). Scoring our model ladder
directly on the Comal Springs series over the 75 rolling origins yields
a performance ranking that mirrors the J-17 head
evaluation:{]}(figs/fig5\_fibre.png)

\textbf{persist:} \(71.9\text{ cfs}\) * \textbf{M1:} \(69.0\text{ cfs}\)
* \textbf{M2m:} \(68.7\text{ cfs}\) * \textbf{M4:} \(73.4\text{ cfs}\) *
\textbf{M3:} \(73.8\text{ cfs}\) * \textbf{M2:} \(74.8\text{ cfs}\) *
\textbf{mean:} \(89.7\text{ cfs}\) * \textbf{M2\_oracle:}
\(45.3\text{ cfs}\)

Crucially, this linear relationship breaks down during severe droughts.
In 1956, when annual-mean J-17 elevation dropped to \(623.15\text{ ft}\)
(with daily levels reaching \(612.51\text{ ft}\)), the linear equation
predicts an annual mean discharge of \(\approx 97\text{ cfs}\), whereas
the spring system actually experienced unprecedented drying (annual mean
of \(32\text{ cfs}\), including complete cessation from June to November
1956).

The zero-discharge intercept of this empirical rating curve occurs at:
\[H_0 = \frac{2876}{4.77} \approx 602.9\text{ ft AMSL}\] which lies well
below any recorded water level at J-17. Because springflow is roughly a
linear transformation of piezometric head under normal conditions,
monitoring spring discharge provides no independent dynamic information
for forecasting J-17, yet simple linear representations fail to capture
the localized non-linearities and threshold transitions that govern
environmental spring drying.

\begin{center}\rule{0.5\linewidth}{0.5pt}\end{center}

\subsubsection{5.6 Pumpage Policy
Counterfactuals}\label{pumpage-policy-counterfactuals}

To evaluate the operational impact of regulatory extraction controls
relative to natural climatic variability, we simulate four
counterfactual pumping scenarios across the post-1990 period
(1991--2023).

Simulations are executed by passing counterfactual pumping paths through
the pre-permitting stock-flow map (fit to 1980--1990 data;
\(\hat{\beta} = 0.0175\),
\(\hat{\gamma} = -0.0284\text{ ft}\cdot(10^3\text{ acre-ft})^{-1}\),
\(\hat{\delta} = -0.254\)) forced with realized historical recharge,
initialized from the observed 1990 head (\(645.8\text{ ft}\); Table 8).

\subsubsection{Table 8. Counterfactual pumping scenarios evaluated
across 1991--2023 using the pre-permit
model.}\label{table-8.-counterfactual-pumping-scenarios-evaluated-across-19912023-using-the-pre-permit-model.}

\emph{Note: Observed terminal 2023 J-17 elevation is
\(635.68\text{ ft}\). The 1980--1990 training model exhibits an OLS
simultaneity bias, which must be considered when interpreting scenario
trajectories.} \textbar{} Counterfactual Pumping Scenario \textbar{}
Mean Annual Pumping (\(10^3\text{ acre-ft}\cdot\text{yr}^{-1}\))
\textbar{} Projected 2023 Head (ft) \textbar{} Simulated Period Minimum
Head (ft) \textbar{} RMSE vs.~Realized Heads (ft) \textbar{} \textbar{}
:--- \textbar{} :---: \textbar{} :---: \textbar{} :---: \textbar{} :---:
\textbar{} \textbar{} \textbf{Realized Historical Pumping} \textbar{}
382.1 \textbar{} 641.6 \textbar{} 640.2 \textbar{} 8.56 \textbar{}
\textbar{} \textbf{Pumping Frozen at 1990 Peak} \textbar{} 489.4
\textbar{} 630.9 \textbar{} 630.5 \textbar{} 14.22 \textbar{} \textbar{}
\textbf{Pre-Permit Mean Pumping} \textbar{} 469.8 \textbar{} 632.4
\textbar{} 632.0 \textbar{} 13.02 \textbar{} \textbar{}
\textbf{Mandatory 20\% Pumping Cut} \textbar{} 305.7 \textbar{} 646.8
\textbar{} 645.5 \textbar{} \textbf{7.19} \textbar{}

These counterfactual simulations illustrate four key dynamics: 1.
\textbf{Pumping Operates as a Secondary Control:} Across the full range
of policy scenarios---from an aggressive 20\% demand reduction up to
freezing extractions at the historical 1990 peak---projected 2023 head
levels span a range of \(15.9\text{ ft}\)
(\(630.9\text{ to } 646.8\text{ ft}\)). This policy spread is of the
same magnitude as the model's own out-of-sample error band
(\(7.19\text{ to } 14.22\text{ ft}\)). 2. \textbf{Recharge Variability
Dominates Piezometric Head:} While the extraction coefficient
(\(\hat{\gamma} = -0.0284\)) is larger in magnitude than the recharge
coefficient (\(\hat{\beta} = 0.0175\)), natural recharge varies over a
vastly wider range
(\(44\text{ to } 2,486 \times 10^3\text{ acre-ft}\cdot\text{yr}^{-1}\),
a span of nearly \(2,440 \times 10^3\text{ acre-ft}\)) than regional
pumping
(\(300\text{ to } 540 \times 10^3\text{ acre-ft}\cdot\text{yr}^{-1}\), a
span of \(240 \times 10^3\text{ acre-ft}\)). An extreme recharge event
of \(1,000 \times 10^3\text{ acre-ft}\) shifts J-17 water levels by
\(17.5\text{ ft}\), dwarfing the marginal impact of annual regulatory
pumping adjustments. 3. \textbf{Apparent Accuracy Gains from Policy
Reductions Reflect Model Bias:} Simulating a 20\% pumping cut yields the
lowest RMSE against observed heads (\(7.19\text{ ft}\)
vs.~\(8.56\text{ ft}\)). Rather than indicating that historical pumping
was suboptimal, this outcome occurs because the pre-permitting model has
an intrinsic positive head bias (projecting a 2023 head of
\(641.6\text{ ft}\) under realized pumping vs.~the observed
\(635.68\text{ ft}\)). Imposing artificial pumping cuts shifts the
baseline trajectory, masking structural over-prediction. 4.
\textbf{Estimated Regulatory Effects are Modest:} A permanent 20\%
extraction reduction increases simulated 2023 water levels by
\(5.2\text{ ft}\) (\(646.8\text{ ft}\) vs.~\(641.6\text{ ft}\); the
actual-pumpage path itself ends \(5.9\text{ ft}\) above the observed
2023 head of \(635.68\text{ ft}\)). While hydrologically meaningful
during drought emergencies, this increment is smaller than the model's
single-year prediction error, demonstrating that regulatory demand
management cannot substitute for predictive skill in operational
planning.

\begin{center}\rule{0.5\linewidth}{0.5pt}\end{center}

The full-sample coefficients
\((\hat{\beta}, \hat{\gamma}) = (0.017, -0.026)\) differ from the
pre-permit window pair because the windows differ; the magnitude
ordering, not the signed pair, carries readings 1--4 above.

\subsection{6. Discussion}\label{discussion}

\subsubsection{6.1 The Timing Bottleneck: Forecasting
vs.~Nowcasting}\label{the-timing-bottleneck-forecasting-vs.-nowcasting}

The primary empirical finding of this study is that dynamic stock-flow
water-balance models forced by persisted fluxes fail to beat naive
persistence for annual-mean groundwater levels. This failure is rooted
in an environmental timing mismatch:

\[\operatorname{corr}(R_t, R_{t-1}) = 0.17 \quad \text{versus} \quad \operatorname{corr}(\Delta H_t, R_t) = 0.74\]

Annual recharge across the Edwards Aquifer recharge zone behaves
essentially as a temporally uncorrelated, near-white noise process. In
contrast, annual water-level increments at J-17 are strongly coupled to
contemporaneous annual recharge.

When a stock-flow model persists antecedent recharge
(\(\tilde{R}_{t+1} = R_t\)), it projects water levels using an
environmental driver that shares less than 3\% of its variance with next
year's actual flux. The model incorporates obsolete boundary forcing,
degrading performance below that of a naive rule that assumes no change.

When realized fluxes are provided exogenously (M2\_oracle), the one-pool
balance achieves an RMSE of \(7.55\text{ ft}\)---a 43\% reduction in
error relative to persistence (\(13.23\text{ ft}\)) and a 53\% reduction
at \(h=5\) (\(10.87\text{ ft}\) vs.~\(21.11\text{ ft}\)).

This contrast defines the operational boundary of lumped water balances
in rapidly responsive aquifers: * Provided with contemporaneous boundary
fluxes, the one-pool water balance functions as an accurate
\textbf{annual nowcast}. * Deployed at a forecast origin with unknown
future fluxes, the identical map fails as an \textbf{operational
forecast}.

This same structural limitation appears in the companion evaluation of
Northern cod (Abaee, 2026). In that marine system, supplying
reconstructed annual fisheries removals fails to improve multi-step
biomass forecasts under a stationary surplus-production model, because
the underlying demographic driver---an abrupt escalation in natural
mortality---is absent from the predictive equation.

In both systems, structural elaboration fails out-of-sample because the
dominant physical driver cannot be reliably projected using information
available at the forecast origin.

\begin{center}\rule{0.5\linewidth}{0.5pt}\end{center}

\subsubsection{6.2 Autocorrelation Mechanics and the Baseline
AR(1)}\label{autocorrelation-mechanics-and-the-baseline-ar1}

The marginal performance of the autonomous AR(1) model (M1), which
improves one-year RMSE over persistence by only \(0.39\text{ ft}\), can
be explained through the time-series properties of index-well water
levels.

An estimated AR(1) process: \[H_{t+1} = a + \varphi H_t\] nests naive
persistence as a parametric special case when \(\varphi = 1\) and
\(a = 0\). As the true autocorrelation of a time series approaches
unity, the estimated autoregression converges toward persistence, and
any theoretical performance advantage is dwarfed by the parameter
estimation error of \(\hat{\varphi}\).

At J-17, the annual head record exhibits moderate persistence, but
remains far from an integrated unit-root process: the full-sample lag-1
autocorrelation is \(\operatorname{corr}(H_t, H_{t-1}) = 0.64\), and
rolling-origin regressions yield a mean coefficient of
\(\hat{\varphi} = 0.66\).

This moderate autocorrelation creates a narrow window where an unforced
AR(1) model can capture slight mean-reverting tendencies without
overfitting. However, because the true series is regularly perturbed by
high-variance recharge shocks, this advantage remains small
(\(0.39\text{ ft}\)) and statistically indistinguishable from zero
(bootstrap CI: {[}\(-1.51, +0.71\){]} ft).

This mechanism leads to a clear, testable hydrogeological hypothesis:
across regional aquifers, as the temporal persistence of the head record
increases (e.g., in systems with thick unconfined vadose zones or high
matrix storage), the performance margin of an AR(1) over naive
persistence should shrink toward zero and eventually turn negative as
estimation noise dominates.

An exploratory test on index well J-27 in the neighboring Uvalde Pool of
the Edwards Aquifer, where annual water levels are more strongly
autocorrelated (\(\operatorname{AC}(1) = 0.84\)), confirms this
expectation: there, the AR(1) model loses to naive persistence by
\(+0.13\text{ ft}\).

Thus, the retention of M1 at J-17 is a consequence of moderate
autocorrelation in a dynamic karst system, and should not be interpreted
as evidence that unforced statistical models are universally superior.

\begin{center}\rule{0.5\linewidth}{0.5pt}\end{center}

\subsubsection{6.3 Benchmark Discipline in Hydrogeological
Forecasting}\label{benchmark-discipline-in-hydrogeological-forecasting}

These results offer an important perspective on the growing groundwater
benchmarking literature. Contemporary benchmarking initiatives, such as
the GEMS-GER machine-learning evaluation (Ohmer et al., 2026) and
regional karst deep-learning comparisons (Zhu et al., 2026), have
advanced the discipline by standardizing large-scale multi-model
evaluations.

However, these frameworks typically evaluate models using metrics such
as the Nash--Sutcliffe Efficiency (\(\mathrm{NSE} > 0.5\)) or
comparative \(R^2\), without testing whether candidate architectures
outperform naive persistence under pre-registered retention rules.

When evaluated under an out-of-sample retention gate, structural
complexity must earn its place by demonstrably outperforming the
baseline prediction that hydrological conditions remain unchanged. In
the Edwards Aquifer, dynamic stock-flow equations, residual tracking,
and antecedent climate predictors all fail this basic test at an annual
horizon.

Furthermore, our multi-step results challenge standard modeling
assumptions: over five-year planning horizons, long-term historical
climatology (\(16.80\text{ ft}\)) decisively outperforms both
persistence (\(21.11\text{ ft}\)) and dynamic water balances
(\(33.49\text{ ft}\)).

In rapidly recharged, regulated karst basins, multi-year water-level
expectations are bounded by environmental mean reversion. Dynamic models
that persist transient drought or pluvial conditions compound error over
time, diverging from observed dynamics.

\begin{center}\rule{0.5\linewidth}{0.5pt}\end{center}

\subsubsection{6.4 Hydrogeological Interpretations of the Lumped
Model}\label{hydrogeological-interpretations-of-the-lumped-model}

The empirical failure of the one-pool balance out-of-sample does not
imply that regional groundwater dynamics are non-physical. Instead, it
reflects the difficulty of capturing complex karst hydrogeology within a
discrete-time scalar mapping: * \textbf{Drainage Mechanics:} The
autonomous AR(1) model M1 can be interpreted physically. Regional spring
discharge from the Edwards Aquifer approximately follows Darcy's law,
scaling linearly with piezometric head above local spring orifices
(\(H_s\)). The analytical solution for unforced aquifer drainage yields
an affine autoregressive relationship:
\[H_{t+1} \approx (1 - k) H_t + k H_s\] where \(k\) represents an annual
aquifer drainage constant. The empirical autoregressive slope
(\(\hat{\varphi} = 0.66\)) corresponds to a drainage decay rate of:
\[\hat{k} = 1 - \hat{\varphi} \approx 0.34\text{ yr}^{-1}\] This
indicates that M1 performs well not because it is an abstract
statistical curve, but because it captures the natural free-drainage
recession of the aquifer when boundary recharge is absent. *
\textbf{Unresolved Karst Geometry:} The one-pool model inevitably
aggregates diverse hydrogeological processes into a single equation.
Localized conduit flow, rapid unconfined recharge through sinking
streams, delayed diffuse infiltration through the unsaturated zone,
cross-fault transfers across the Balcones Fault Zone, and pressure-wave
propagation across the confined zone are all relegated to the residual
term \(\varepsilon_t\).

While distributed numerical models (e.g., MODFLOW) attempt to resolve
this spatial heterogeneity directly, they require extensive boundary
inputs and calibration. Our results demonstrate that simplified, lumped
conceptual representations cannot overcome the timing bottleneck:
without reliable predictions of future boundary fluxes, adding lumped
physical equations fails to improve out-of-sample forecasts.

\begin{center}\rule{0.5\linewidth}{0.5pt}\end{center}

\subsubsection{6.5 Methodological
Limitations}\label{methodological-limitations}

To support balanced application of these conclusions, several
limitations of the evaluation design must be highlighted: 1.
\textbf{Flux Aggregation Across Sub-Pools:} Reconstructed recharge
(\(R_t\)) and extraction (\(P_t\)) represent totals across the wider San
Antonio and Uvalde basins. While J-17 serves as the index well for the
San Antonio Pool, inter-pool transfers across the Knippa Gap divide
introduce unmodeled boundary fluxes that can distort local mass
balances. 2. \textbf{Estimation Uncertainty in Historical Fluxes:}
Recharge estimates derived from the Puente stream-loss method are
subject to unquantified streamflow gauging and catchment runoff errors.
Similarly, reported well extractions exclude unmetered domestic,
livestock, and historical federal withdrawals. 3. \textbf{Temporal
Resolution Constraints:} This evaluation operated strictly at an annual
time step. While well suited for annual water-budget planning and
multi-year drought management, this temporal resolution aggregates
high-frequency recharge pulses and localized summer drawdown spikes.
Testing whether sub-annual (e.g., monthly or seasonal) water balances
demonstrate skill requires an independent, sub-annual evaluation
protocol. 4. \textbf{Information Latency vs.~Operational Monitoring:}
Module M4 evaluated a theoretical one-year lag in initial head
observations, matching the structure of our companion marine assessment.
In practice, J-17 is a modern telemetered monitoring well with real-time
daily data availability. Consequently, M4 quantifies the mathematical
cost of stale initial conditions rather than a real-world constraint of
the Edwards monitoring network.

\begin{center}\rule{0.5\linewidth}{0.5pt}\end{center}

\subsection{7. Conclusions}\label{conclusions}

Evaluating a ladder of discrete-time water-balance models against 90
years of annual water-level observations at index well J-17 yields four
main conclusions: 1. \textbf{Dynamic Stock-Flow Models Fail to Beat
Persistence:} A lumped one-pool water-balance model that persists
prior-year recharge fails to outperform naive last-value persistence at
a one-year forecast horizon (\(14.70\text{ ft}\)
vs.~\(13.23\text{ ft}\)). Annual recharge across the Edwards Aquifer
behaves as near-white noise
(\(\operatorname{corr}(R_t, R_{t-1}) = 0.17\)), so forcing forward
predictions with persisted fluxes introduces substantial error into the
state equation. 2. \textbf{Univariate Autoregression Provides a
Marginal, Statistically Fragile Baseline:} An unforced linear AR(1)
model achieves a minor one-year RMSE reduction of \(0.39\text{ ft}\)
(\(12.84\text{ ft}\)), which clears the frozen point-retention rule.
However, this margin is statistically indistinguishable from zero (mean
absolute error is tied at \(10.7\text{ ft}\), and the 95\% moving-block
bootstrap confidence interval spans zero: {[}\(-1.51, +0.71\){]} ft).
This performance captures moderate empirical autocorrelation
(\(\hat{\varphi} \approx 0.66\)) reflecting natural baseflow recession,
rather than an underlying stock-flow dynamic. 3. \textbf{Water Balances
Function as Nowcasts Rather than Operational Forecasts:} Provided with
realized boundary fluxes, the one-pool balance achieves an RMSE of
\(7.55\text{ ft}\) (a 43\% error reduction over persistence). At an
annual time step, the lumped model functions effectively as an
environmental nowcast. However, antecedent climate indices (ENSO SST
anomalies and division rainfall) improve one-year head RMSE by at most
\(0.13\text{ ft}\), failing to bridge the gap between nowcasting and
operational forecasting. 4. \textbf{Historical Climatology is the Robust
Benchmark for Multi-Year Planning:} Over five-year horizons, long-term
training-mean climatology decisively outperforms both naive persistence
(\(16.80\text{ ft}\) vs.~\(21.11\text{ ft}\); bootstrap interval
strictly excludes zero) and dynamic models forced by persisted fluxes
(\(33.49\text{ ft}\)). In rapidly recharged karst aquifers, multi-year
expectations are bounded by environmental mean reversion.

\subsubsection{Practical Recommendations for Groundwater
Managers}\label{practical-recommendations-for-groundwater-managers}

\begin{itemize}
\tightlist
\item
  \textbf{Exercise Caution with Persisted-Flux Predictions:} At annual
  decision origins, water managers should avoid using lumped
  water-balance forecasts forced by persisted hydrologic conditions.
  Without reliable forward predictions of regional recharge, these
  models perform worse than simple last-value persistence.
\item
  \textbf{Adopt Climatological Defaults for Multi-Year Planning:} For
  medium-term (3- to 5-year) drought-mitigation and water-supply
  planning, managers should rely on long-term historical climatological
  baselines rather than forward projections from dynamic balances forced
  by recent dry conditions, which tend to project unrealistically
  severe, persistent drawdowns.
\item
  \textbf{Recognize the Limits of Pumping Adjustments:} While regulatory
  extraction controls remain critical for mitigating extreme drawdowns
  during severe droughts, counterfactual simulations demonstrate that
  regional pumping is a secondary lever compared to natural recharge
  variability. Expected water-level shifts from demand management fall
  within the predictive uncertainty of annual models.
\end{itemize}

\subsubsection{Scope and
Transferability}\label{scope-and-transferability}

These findings apply directly to shallow, rapidly recharged,
institutionally managed regional aquifers characterized by high
hydrologic conductivity and responsive boundary fluxes. They do not
apply to deep, non-stationary fossil aquifer systems undergoing
continuous, unreplenished depletion, where water levels follow monotonic
trends and naive persistence fails.

Ultimately, this study demonstrates the value of benchmark discipline:
before adopting complex predictive models for operational water
resources management, their structural additions must be quantitatively
justified against simple persistence and climatological baselines using
locked, out-of-sample evaluation protocols.

\begin{center}\rule{0.5\linewidth}{0.5pt}\end{center}

\subsection{Data and Code
Availability}\label{data-and-code-availability}

All hydrologic observation series, protocol documentation, model
execution scripts, and statistical replication codes are open-source and
archived in public repositories to ensure computational reproducibility:
1. \textbf{Primary Observation Records:} J-17 daily water levels are
available from the Texas Water Development Board Groundwater Database
(Well 68-37-203). Annual regional recharge and pumping compilations are
maintained by the Edwards Aquifer Authority and the U.S. Geological
Survey. Climate teleconnection indices are archived by NOAA CPC (Niño
3.4) and NOAA NCEI (nClimDiv). 2. \textbf{Computational Scripts:} The
complete evaluation pipeline, parameter estimation routines,
Diebold--Mariano tests, and moving-block bootstrap routines are
permanently archived under repository tag \texttt{wave\_e\_edwards}
(with independent replication code in
\texttt{campaign\_e3\_dm\_uncertainty.py}).

\begin{center}\rule{0.5\linewidth}{0.5pt}\end{center}

\subsection{References}\label{references}

\begin{itemize}
\item
  Adamowski, J., and Chan, H.F. 2011. A wavelet neural network
  conjunction model for groundwater level forecasting. \emph{Journal of
  Hydrology} 407: 28--40. https://doi.org/10.1016/j.jhydrol.2011.06.013
\item
  Abaee, A. 2026. Does a surplus-production ladder improve forecasts of
  Northern cod? A scored test on NAFO 2J3KL. Zenodo.
  https://doi.org/10.5281/zenodo.22553609. Companion forecast-evaluation
  study (Northern cod, NAFO 2J3KL).
\item
  Daliakopoulos, I.N., Coulibaly, P., and Tsanis, I.K. 2005. Groundwater
  level forecasting using artificial neural networks. \emph{Journal of
  Hydrology} 309: 229--240.
  https://doi.org/10.1016/j.jhydrol.2004.12.001
\item
  Diebold, F.X., and Mariano, R.S. 1995. Comparing predictive accuracy.
  \emph{Journal of Business \& Economic Statistics} 13: 253--263.
  https://doi.org/10.1080/07350015.1995.10524599
\item
  Edwards Aquifer Authority. 2024/25. \emph{2023 Groundwater Discharge
  and Usage}, Table 1 (after USGS letter report, 5 April 2024).
  https://www.edwardsaquifer.org/wp-content/uploads/2025/05/2023-Groundwater-Discharge-and-Usage.pdf
\item
  Edwards Aquifer Authority. Critical Period / Drought Management.
  https://www.edwardsaquifer.org/groundwater-users/critical-period-drought-management/
\item
  Künsch, H.R. 1989. The jackknife and the bootstrap for general
  stationary observations. \emph{Annals of Statistics} 17: 1217--1241.
  https://doi.org/10.1214/aos/1176347265
\item
  Makridakis, S., Spiliotis, E., and Assimakopoulos, V. 2020. The M4
  Competition: 100,000 time series and 61 forecasting methods.
  \emph{International Journal of Forecasting} 36: 54--74.
\item
  NOAA National Centers for Environmental Information. nClimDiv
  precipitation (climdiv-pcpndv-v1.0.0).
  https://www.ncei.noaa.gov/pub/data/cirs/climdiv/
\item
  NOAA Physical Sciences Laboratory. Niño 3.4 monthly SST (HadISST).
  https://psl.noaa.gov/data/timeseries/month/data/nino34.long.data
\item
  Ohmer, M., Liesch, T., Habbel, B., Heudorfer, B., Gomez, M., Clos, P.,
  Nölscher, M., and Broda, S. 2026. GEMS-GER: A machine learning
  benchmark dataset of long-term groundwater levels in Germany with
  meteorological forcings and site-specific environmental features.
  \emph{Earth System Science Data} 18, 77.
\item
  Puente, C. 1978. \emph{Method of estimating natural recharge to the
  Edwards Aquifer in the San Antonio area, Texas.} U.S. Geological
  Survey, Austin, Texas.
\item
  Ropelewski, C.F., and Halpert, M.S. 1986. North American precipitation
  and temperature patterns associated with the El Niño/Southern
  Oscillation (ENSO). \emph{Monthly Weather Review} 114: 2352--2362.
\item
  Scanlon, B.R., Mace, R.E., Barrett, M.E., and Smith, B. 2003. Can we
  simulate regional groundwater flow in a karst system using equivalent
  porous media models? Case study, Barton Springs Edwards aquifer, USA.
  \emph{Journal of Hydrology} 276: 137--158.
\item
  Texas Water Development Board. Water Data for Texas, well 6837203
  (J-17). https://waterdatafortexas.org/groundwater/well/6837203
\item
  Umphres, G.D., and Choi, N.J. 2025. Estimated Annual Recharge to the
  Edwards Aquifer in the San Antonio Area, by Stream Basin or Ungaged
  Area, 1934--2024. U.S. Geological Survey data release.
  https://doi.org/10.5066/P1BI62NY
\item
  U.S. Geological Survey. National Water Information System, site
  08168710, Comal Springs at New Braunfels, Texas.
\item
  Zhu, Q., Zhu, Y., Niu, J., Huang, J., Huang, F., Zhou, X., Liu, D.,
  and Hu, B.X. 2026. Benchmarking machine learning and deep learning
  models for groundwater level prediction in karst aquifers: The
  dominant role of hydrogeological complexity. \emph{Water} 18, no. 8:
  939.
\end{itemize}

\subsection{Declarations}\label{declarations}

\subsubsection{Data availability}\label{data-availability}

All input data, analysis scripts, and result files are archived in the
public repository at
https://github.com/MIKEAA2020/general-sustainability, together with the
frozen scoring protocols (dated 2026-08-25, locked before any score was
generated). J-17 daily highs: Texas Water Development Board well
6837203. Recharge: Umphres and Choi (2025), USGS data release,
https://doi.org/10.5066/P1BI62NY. Pumpage: Edwards Aquifer Authority
Table 1. Comal Springs: USGS 08168710. Niño 3.4: NOAA PSL HadISST (raw
file registered with the repository). Precipitation: NCEI nClimDiv ---
the raw file is not distributed with the repository (provenance URL
archived in the sources index), so the three precipitation columns of
the fixed panel are not reproducible from the registered code alone,
while the two Niño columns rebuild from the registered file to machine
precision; scoring from the registered analysis panel does not require
the nClimDiv file. The registered twenty-column analysis panel is the
dataset of record for all scored analyses. All computations are
deterministic: re-executing the registered scripts in a fresh
environment regenerated every archived result file byte for byte, and
all scored rows recompute from the per-observation forecast files and
the registered series. The pumpage counterfactual layer (Section 5.6) is
produced by
\texttt{rerun\_campaigns/campaign\_e3\_pumpage\_scenarios.py}, archived
alongside its outputs, and regenerates them exactly. The post-freeze
uncertainty layer (Section 5.3.1) is produced by
\texttt{wave\_e\_edwards/src/e3\_audit\_uncertainty.py} (seeded,
deterministic; Diebold--Mariano with Newey--West HAC and moving-block
bootstrap on the archived per-origin forecast files), with its outputs
archived as
\texttt{wave\_e\_edwards/results/e3\_audit\_uncertainty.json}; the
clip-binding statement of Section 5.2 is computed there from the
registered panel and reproduces both fixed-window M2 RMSEs (18.11 and
55.32 ft) exactly. The independent replication of the post-freeze
uncertainty layer is registered as
\texttt{batch\ 7\ (audits\ of\ agent\ arena\ 1\ paper\ rewrites)/campaign\_e3\_dm\_uncertainty.py}
(archived), with its deterministically reproduced output archived
alongside it at
\texttt{batch\ 7\ (audits\ of\ agent\ arena\ 1\ paper\ rewrites)/results/e3\_dm\_uncertainty.csv};
it recomputes the Diebold--Mariano and moving-block-bootstrap table from
the same registered per-origin forecast file under different seeds,
block lengths, and HAC conventions, and every load-bearing conclusion of
Section 5.3.1 is unchanged under either implementation.

\subsubsection{Declaration of competing
interests}\label{declaration-of-competing-interests}

None declared.

\subsubsection{Funding}\label{funding}

No funding was received for this work.

\subsubsection{CRediT authorship contribution
statement}\label{credit-authorship-contribution-statement}

Amin Abaee: conceptualization, data curation, formal analysis,
investigation, methodology, software, validation, visualization, writing
--- original draft, writing --- review and editing.

\subsubsection{Statement on the use of generative artificial
intelligence}\label{statement-on-the-use-of-generative-artificial-intelligence}

Generative-AI tools were used for language editing, register smoothing,
and token-level drift checks of this manuscript. All scientific
statements, numbers, protocol events, and conclusions are the author's;
all AI-produced phrasing was verified against the frozen sources and
archives before release.

\end{document}
